% 数学ノート集・共通テンプレート v2（積分論2026年度第1部の版面を参考）
% A5判。章内で定義・定理・例・注意等の番号を共通化。
% コンパイル： platex -interaction=nonstopmode math_note_template_v2.tex (2回)
%              dvipdfmx math_note_template_v2.dvi
\documentclass[a5paper,10pt,onecolumn]{jsarticle}
\usepackage[dvipdfmx,top=22mm,bottom=20mm,inner=19mm,outer=18mm,headheight=13pt,headsep=6mm]{geometry}
\usepackage{amsmath,amssymb,amsthm,mathtools}
\usepackage{enumitem,fancyhdr,titlesec}
\usepackage[dvipdfmx,bookmarks=true,colorlinks=false]{hyperref}
\usepackage{pxjahyper}

% ----- 差し替える文書情報 -----
\newcommand{\NoteTitle}{Liouville可積分性と微分Galois理論}
\newcommand{\NoteSeries}{数学ノート集}

% ----- 文字・組版 -----
% 参考PDFの質素な明朝＋Computer Modernの印象を引き継ぎながら，
% A5判で窮屈にならないよう本文の行間を少し確保する．
\linespread{1.13}
\setlength{\parindent}{1zw}
\setlength{\parskip}{0pt}
\setlength{\abovedisplayskip}{6pt plus 2pt minus 2pt}
\setlength{\belowdisplayskip}{6pt plus 2pt minus 2pt}
\setlength{\abovedisplayshortskip}{3pt plus 2pt minus 1pt}
\setlength{\belowdisplayshortskip}{4pt plus 2pt minus 1pt}
\setlist[enumerate]{topsep=3pt,itemsep=1pt,parsep=0pt,leftmargin=2.4zw}
\titleformat{\section}{\normalfont\normalsize\bfseries}{\thesection}{0.8zw}{}
\titlespacing*{\section}{0pt}{2.0ex plus .4ex}{.9ex}
\titleformat{\subsection}{\normalfont\normalsize\bfseries}{\thesubsection}{0.7zw}{}
\titlespacing*{\subsection}{0pt}{1.5ex plus .3ex}{.6ex}
\raggedbottom

% ----- 柱とページ番号：ページ上部、罫線なし -----
\fancypagestyle{mathnote}{%
  \fancyhf{}%
  \fancyhead[L]{\footnotesize \NoteSeries\ （\NoteTitle）}%
  \fancyhead[R]{\footnotesize \thepage}%
  \renewcommand{\headrulewidth}{0pt}%
  \renewcommand{\footrulewidth}{0pt}%
}
\pagestyle{mathnote}

% ----- 節内通し番号の定義・定理群 -----
\newtheoremstyle{mathnotetheorem}
 {4.5pt}{4.5pt}{\normalfont}{}%
 {\bfseries}{}{0.65zw}%
 {\thmname{#1}\thmnumber{ #2}\thmnote{（#3）}}
\theoremstyle{mathnotetheorem}
\newtheorem{theorem}{定理}[section]
\newtheorem{definition}[theorem]{定義}
\newtheorem{lemma}[theorem]{補題}
\newtheorem{proposition}[theorem]{命題}
\newtheorem{corollary}[theorem]{系}
\newtheorem{example}[theorem]{例}
\newtheorem{remark}[theorem]{注意}
\renewcommand{\proofname}{証明}
\renewcommand{\qedsymbol}{$\square$}
\makeatletter
\renewenvironment{proof}[1][\proofname]{%
  \par\pushQED{\qed}\normalfont\topsep4\p@\@plus2\p@\relax
  \trivlist\item[\hskip\labelsep\bfseries #1]\ignorespaces
}{\popQED\endtrivlist\@endpefalse}
\makeatother

% ----- 使うノートで適宜拡張する記号 -----
\newcommand{\R}{\mathbb{R}}
\newcommand{\norm}[1]{\lVert #1\rVert}
\newcommand{\abs}[1]{\lvert #1\rvert}

\begin{document}
\section*{読了後にできること}
本稿を読むための前提は，微分体，初等関数，微分Galois群．
本稿の具体的な到達課題は，単純な微分方程式のLiouvillian解を構成し，その条件を確認する．

\section{Liouvillian拡大の定義}

本稿では標数0の微分体 $(K,\partial)$ と
代数閉な定数体 $k=K^\partial$ を考える．
Picard--Vessiot拡大の存在・一意性・Galois対応，
代数群の連結成分・Lie--Kolchinの定理，
代数関数のPuiseux展開を前提とする．

\begin{definition}
微分体拡大 $E/F$ が次のいずれかに属するとき，
一段のLiouvillian拡大という．
\begin{enumerate}
\item $E/F$ は有限代数拡大である．
\item $E=F(t)$ で $\partial t\in F$ を満たす．
\item $E=F(t)$ で $t\ne0$ かつ
 $\partial t/t\in F$ を満たす．
\end{enumerate}
有限の塔
\[
 K=F_0\subset F_1\subset\cdots\subset F_m=E
\]
の各段階がこれらのいずれかであり，さらに各段階で
定数体を増やさない（$F_i^\partial=k$，$0\le i\le m$）とき，
$E/K$ を本稿ではLiouvillian拡大という．
定数体の保存は，後述するPicard--VessiotのGalois対応と
Kolchinの判定定理を適用するための仮定である．
\end{definition}

\begin{definition}
線形微分方程式 $\partial Y=AY$ が
\emph{Liouvillianな解を持つ}とは，
その基本解行列の全成分が
あるLiouvillian拡大に含まれることをいう．
\end{definition}

\begin{remark}
Liouvillian拡大の第三の段階は指数関数を，
第二の段階は不定積分を加える操作に対応する．
これは初等関数だけで表せるという条件より広い．
例えば $e^{-x^2}$ の不定積分は
Liouvillianな拡大として許される．
\end{remark}

\section{一次方程式と微分Galois群}

\begin{proposition}[指数型の方程式]\label{prop:exp}
$a\in F$ とし
\[
 y'=ay
\]
を考える．$t\ne0$ を解として定数体を増やさない
微分体拡大 $E=F(t)$ を取ると，
微分自己同型は
\[
 \sigma(t)=c_\sigma t,\qquad c_\sigma\in k^\times
\]
と作用する．従ってGalois群は
$\mathbb G_m$ の代数部分群である．
\end{proposition}
\begin{proof}
$\partial(\sigma(t)/t)=0$ なので，
$\sigma(t)/t$ は定数となる．
定数体が $k$ のままであるから
$c_\sigma\in k^\times$ である．
Picard--Vessiot理論により群は
$\mathbb G_m$ の閉代数部分群となる．
\end{proof}

\begin{proposition}[積分型の方程式]\label{prop:int}
$a\in F$ とし
\[
 t'=a
\]
を満たす $t$ を加えた
定数体を増やさない微分体拡大 $E=F(t)$ を取る．
微分自己同型は
\[
 \sigma(t)=t+c_\sigma,\qquad c_\sigma\in k
\]
と作用する．従ってGalois群は
$\mathbb G_a$ の代数部分群である．
\end{proposition}
\begin{proof}
$\partial(\sigma(t)-t)=\sigma(a)-a=0$ なので
$\sigma(t)-t$ は定数であり，
定数体は $k$ なので $c_\sigma\in k$ となる．
実際，$t'=a$ に対し
\[
 A=\begin{pmatrix}0&a\\0&0\end{pmatrix},\qquad
 Y=\begin{pmatrix}1&t\\0&1\end{pmatrix}
\]
と置けば $Y'=AY$ である．$E=F(t)$ は定数体を増やさず，
$t$ が $F$ 上代数的なら，最小多項式を微分して
その最高次係数の次の係数を比較すると，
$\partial(t+b)=0$ となる $b\in F$ が得られる．
定数体が $k\subset F$ のままなので $t\in F$ であり，
拡大は自明である．
$t$ が超越的なら，$F[t]$ が微分単純であることを示す．
真の非零微分イデアルがあれば，
そこに属する次数最小のモニック多項式
$P=t^d+b_{d-1}t^{d-1}+\cdots$ を取れる（$d\ge1$）．
$\partial P$ は次数が $d$ より小さく同じイデアルに属するから
$\partial P=0$ となる．$t^{d-1}$ の係数を比較すれば
$da+\partial b_{d-1}=0$，従って
$\partial(t+b_{d-1}/d)=0$ を得る．
これが定数体の保存に反するため，$F[t]$ は微分単純である．
ゆえにいずれの場合も対応するPicard--Vessiot拡大を与える．
従って微分Galois群は
$\left\{\begin{psmallmatrix}1&c\\0&1\end{psmallmatrix}:c\in k\right\}$
の代数部分群となる．
\end{proof}

\section{可解代数群とLiouvillian拡大}

\begin{definition}
線形代数群 $G$ の単位元を含むZariski連結成分を
$G^\circ$ と書く．
群 $G$ が可解であるとは，導来列
$G^{(0)}=G$，
$G^{(i+1)}=[G^{(i)},G^{(i)}]$ が有限回で
自明群となることをいう．
代数群については各導来部分群のZariski閉包を取る．
\end{definition}

\begin{lemma}[連結可解群の一次元商]\label{lem:solvable-series}
$k$ を標数0の代数閉体とし，
$G$ を連結可解線形代数群とする．
このとき正規閉部分群の列
\[
 G=G_0\triangleright G_1\triangleright\cdots
 \triangleright G_r=1
\]
で，各商 $G_i/G_{i+1}$ が
$\mathbb G_a$ または $\mathbb G_m$ に同型となるものが存在する．
\end{lemma}
\begin{proof}
Lie--Kolchinの定理によって，
$G$ は適当な基底で上三角行列群 $B_n$ に含まれる．
上三角群から対角群への準同型
$B_n\to(\mathbb G_m)^n$ の核は単位上三角群 $U_n$ である．
$G$ の像は連結な代数トーラスであり，
適当な基底を選べば
$\mathbb G_m$ を一つずつ取り除く
部分トーラスの正規列を持つ．
また $G\cap U_n$ は連結な冪単代数群である．
標数0の冪単群の指数・対数対応により，
そのLie代数の中心列を群へ戻すと，
一次元の加法群 $\mathbb G_a$ を商とする
正規列を得る．
これらの正規列を $G$ 内へ引き戻して連結すると，
各商が $\mathbb G_a$ または $\mathbb G_m$ となる
所望の列を得る．
\end{proof}

\begin{lemma}[加法・乗法群のtorsorの自明性]
\label{lem:hilbert90}
標数0の体 $F$ 上の $\mathbb G_a$-torsorおよび
$\mathbb G_m$-torsorは自明である．
\end{lemma}
\begin{proof}
まず有限Galois拡大 $E/F$ 上でtorsorを自明化し，
Galois群 $\Gamma$ に関する降下データを考える．
加法群の場合，降下データは
\[
 c_{\sigma\tau}=c_\sigma+\sigma(c_\tau)
\]
を満たす元 $c_\sigma\in E$ の族として表される．
\[
 b=-\frac1{|\Gamma|}\sum_{\tau\in\Gamma}c_\tau
\]
と置くと，
\[
 \sigma(b)-b=c_\sigma
\]
がコサイクルの等式から従う．
従って座標を $t\mapsto t-b$ と変更すれば
降下データが消え，$F$ 上のtorsorが自明化する．

乗法群の場合には
$c_{\sigma\tau}=c_\sigma\sigma(c_\tau)$ とする．
適当な $a\in E$ を選び
\[
 b=\sum_{\tau\in\Gamma}c_\tau\tau(a)\ne0
\]
となるようにできる．
これは異なる体自己同型が $E$ 上線形独立であるという
Artinの独立性定理による．
コサイクルの等式から
$\sigma(b)=c_\sigma^{-1}b$ となるので，
$c_\sigma=b/\sigma(b)$ である．
従って座標を $b$ で掛け替えることで
乗法的な降下データが消える．
いずれの場合もtorsorは $F$-点を持ち，
自明である．
\end{proof}

\begin{proposition}[$\mathbb G_a$ および $\mathbb G_m$ 型PV拡大]
\label{prop:one-dim-pv}
Picard--Vessiot拡大 $E/F$ の微分Galois群が
$\mathbb G_a$ に同型なら，
$E=F(t)$，$\partial t\in F$ となる $t$ が存在する．
群が $\mathbb G_m$ に同型なら，
$E=F(t)$，$\partial t/t\in F$ となる
$t\ne0$ が存在する．
\end{proposition}
\begin{proof}
Picard--Vessiot環のtorsor同型から，
$E/F$ は対応する群のtorsorの関数体である．
補題\ref{lem:hilbert90}によってこのtorsorは
$F$ 上自明である．
$\mathbb G_a$ の場合は座標 $t$ を選ぶと，
群作用は $\sigma(t)=t+c_\sigma$ となる．
導分が群作用と可換なので
$\sigma(\partial t)=\partial t$ である．
微分Galois対応から固定体は $F$ なので
$\partial t\in F$ となる．
torsorの関数体は $F(t)$ だから $E=F(t)$ である．
$\mathbb G_m$ の場合も
$\sigma(t)=c_\sigma t$ となる座標を取れる．
すると $\sigma(\partial t/t)=\partial t/t$ なので
$\partial t/t\in F$ であり，
同様に $E=F(t)$ となる．
\end{proof}

\section{Kolchinの可積分性判定}

\begin{theorem}[Kolchinの定理]\label{thm:kolchin}
$L/K$ を線形方程式 $\partial Y=AY$ の
Picard--Vessiot拡大とし，
$G=\operatorname{Gal}^\partial(L/K)$ とする．
次は同値である．
\begin{enumerate}
\item 方程式はLiouvillian拡大の中で基本解行列を持つ．
\item $L/K$ はLiouvillian拡大である．
\item $G^\circ$ は可解代数群である．
\end{enumerate}
\end{theorem}
\begin{proof}
(2) なら基本解行列は $L$ の中にあるので (1) は明らかである．
まず (3) から (2) を示す．
$G^\circ$ は有限指数の正規閉部分群なので，
微分Galois対応から
$F=L^{G^\circ}$ は $K$ の有限代数拡大である．
また $L/F$ はPicard--Vessiot拡大で，
微分Galois群は $G^\circ$ に等しい．
補題\ref{lem:solvable-series}によって
\[
 G^\circ=G_0\triangleright\cdots\triangleright G_r=1
\]
という正規列を取る．
$F_i=L^{G_i}$ と置くとGalois対応から
\[
 F=F_0\subset F_1\subset\cdots\subset F_r=L
\]
となる．
各 $F_{i+1}/F_i$ はPicard--Vessiot拡大であり，
そのGalois群は $G_i/G_{i+1}$ である．
命題\ref{prop:one-dim-pv}により，
各段階は積分型または指数積分型の一段拡大となる．
最初に有限代数拡大 $F/K$ を付け加えれば
$L/K$ はLiouvillian拡大である．

逆方向を示す．
(1) により基本解を含むLiouvillian拡大 $E/K$ が存在する．
$L$ は基本解行列とその逆行列の成分で生成され，
定数体を増やさない最小の微分体なので，
共通の微分閉包の中で $L\subset E$ とみなせる．
Liouvillian塔の各段階を正規化すると，
代数段階は有限Galois拡大，
積分段階と指数積分段階はそれぞれ
$\mathbb G_a$，$\mathbb G_m$ をGalois群に持つ
Picard--Vessiot拡大として扱える．
この正規化で得られる拡大のGalois群は，
有限群，$\mathbb G_a$，$\mathbb G_m$ を商とする
有限個の群拡大として記述される．
有限群の単位成分は自明であり，
$\mathbb G_a$ と $\mathbb G_m$ は可換なので，
こうした群の単位成分は可解である．
$L$ のGalois群はそのような群の閉部分群の商として得られるため，
$G^\circ$ も可解である．
従って (1) から (3) を得る．
\end{proof}

\begin{remark}
定理\ref{thm:kolchin}の逆方向では
Liouvillian塔の正規化に関する微分代数的降下を用いた．
この降下は，通常の体のGalois閉包に相当する構成と，
微分体における定数体の保存を組み合わせるものである．
\end{remark}

\section{二階線形方程式とRiccati方程式}

\begin{definition}
$K$ 上の二階線形方程式
\[
 y''=r y,\qquad r\in K
 \tag{R}
\]
を考える．
\[
 u'+u^2=r
 \tag{Ric}
\]
を対応するRiccati方程式という．
\end{definition}

\begin{lemma}[対数微分とRiccati方程式]
\label{lem:riccati-reduction}
$y\ne0$ が (R) の解なら $u=y'/y$ は (Ric) の解である．
逆に (Ric) の解 $u$ に対して
$y'=uy$ の非零解 $y$ は (R) の解となる．
\end{lemma}
\begin{proof}
$u=y'/y$ なら
\[
 u'=\frac{y''y-(y')^2}{y^2}
 =\frac{y''}{y}-u^2,
\]
従って $u'+u^2=y''/y=r$ である．
逆に $y'=uy$ を微分すると
$y''=(u'+u^2)y=ry$ となる．
\end{proof}

\begin{lemma}[第二解の積分表示]\label{lem:second-solution}
$y$ を (R) の非零解とすると，
線形独立なもう一つの解は
\[
 z=y\int y^{-2}\,dx
\]
で与えられる．
\end{lemma}
\begin{proof}
$v' =y^{-2}$ となる $v$ を加え，$z=yv$ と置く．
\[
 z''=y''v+2y'v'+yv''
\]
であり，
$v''=-2y^{-3}y'$ なので後二項は相殺する．
従って $z''=y''v=ryv=rz$ である．
Wronski行列式は
\[
 yz'-y'z=y(y'v+yv')-y'(yv)=y^2v'=1
\]
なので二解は線形独立である．
\end{proof}

\begin{theorem}[二階方程式の可積分性判定]
\label{thm:riccati-criterion}
$k$ を代数閉とし，$L/K$ を (R) のPicard--Vessiot拡大，
$G$ を微分Galois群とする．
次は同値である．
\begin{enumerate}
\item $G^\circ$ は可解である．
\item (Ric) は $K$ 上代数的な微分体拡大の中に解を持つ．
\item (R) はLiouvillian拡大の中で基本解を持つ．
\end{enumerate}
\end{theorem}
\begin{proof}
(1) ならLie--Kolchinの定理によって
$G^\circ$ は二次元解空間の一次元部分空間を保つ．
従って非零解 $y$ を選んで
$\sigma(y)=c_\sigma y$ が全ての $\sigma\in G^\circ$
で成立するようにできる．
このとき $u=y'/y$ は
全ての $\sigma\in G^\circ$ で不変なので
$u\in L^{G^\circ}$ である．
$L^{G^\circ}/K$ は有限代数拡大だから
$u$ は $K$ 上代数的であり，
補題\ref{lem:riccati-reduction}から (Ric) を満たす．
よって (2) が従う．

(2) とし，$u$ を代数拡大 $K(u)$ に取る．
まず $y'=uy$ の非零解 $y$ を指数積分型に加えれば，
補題\ref{lem:riccati-reduction}から (R) の解になる．
続いて $v'=y^{-2}$ を積分型に加えると，
補題\ref{lem:second-solution}から $yv$ が第二解となり，
二解のWronski行列式は1なので基本解系をなす．
従って方程式はLiouvillianな基本解を持ち (3) が従う．
(3) から (1) は定理\ref{thm:kolchin}による．
\end{proof}

\section{計算可能な具体例}

\begin{example}[指数型の基本解]
$K=\mathbb C(x)$，$r(x)=1$ とする．
$y''=y$ の二つの基本解は $e^x,e^{-x}$ であり，
Liouvillianである．
\end{example}
\begin{proof}
$e^x$，$e^{-x}$ は $y''=y$ を満たし，
Wronski行列式は $-2$ なので線形独立である．
$e^x$ は $y'/y=1$ を満たす指数積分型の元であり，
その逆元も同じ微分体に属する．
従って基本解系はLiouvillian拡大に含まれる．
\end{proof}

\begin{example}[積分が必要な第二解]
$K=\mathbb C(x)$ とし，
\[
 y''=\left(1+\frac2{x^2}\right)y
\]
を考える．
$y_1=e^x(1-x^{-1})$ は非零解であり，
第二解は $y_1\int y_1^{-2}\,dx$ として得られる．
\end{example}
\begin{proof}
$y_1'=e^x(1-x^{-1}+x^{-2})$，
$y_1''=e^x(1-x^{-1}+2x^{-2}-2x^{-3})$ なので
\[
 y_1''=\left(1+\frac2{x^2}\right)e^x(1-x^{-1})
\]
が成立する．
補題\ref{lem:second-solution}によって第二解が得られ，
これに必要な操作は指数積分と不定積分だけである．
\end{proof}

\section{Airy方程式の非Liouvillian性}

\begin{lemma}[代数関数のRiccati解に関する補題]
\label{lem:airy-riccati}
微分体 $K=\mathbb C(x)$ において
\[
 u'+u^2=x
 \tag{A-Ric}
\]
は $K$ 上代数的な解を持たない．
\end{lemma}
\begin{proof}
背理法で，$u$ を $\mathbb C(x)$ 上代数的な解とする．
代数関数のNewton--Puiseuxの定理により，
無限遠点のある分枝で $u$ は
$x^{-1/m}$ の収束するPuiseux級数に展開できる．
$x$ の最高次を比較すると，
$u^2$ の最高次が $x$ と一致する必要がある．
実際微分は最高次数を1だけ下げるので，
最高次数を $\alpha$ とすると
$\max\{\alpha-1,2\alpha\}=1$ であり，
これを満たすのは $\alpha=1/2$ である．
従って $u\sim\varepsilon x^{1/2}$，
$\varepsilon\in\{1,-1\}$ となる．

$z=x^{-3/2}$ と置き，
$u=x^{1/2}v(z)$ と表す．
形式的な係数比較によって
\[
 v(z)^2+z\left(\frac12v(z)-\frac32zv'(z)\right)=1
 \tag{7}
\]
を得る．
$\varepsilon$ を指定した形式的Puiseux解は，
この等式によって逐次一意に決まる．
実際，先頭の次数を決めると，
新しく現れる係数は $2\varepsilon$ 倍されて現れ，
導関数による項はより低次の係数から決まるからである．
従って解に任意の別の分数指数の項を加える余地はなく，
\[
 v(z)=\sum_{n=0}^\infty c_nz^n,\qquad c_0=\varepsilon
\]
という形式級数が唯一の候補となる．
式 (7) の $z^n$ の係数を比較すると
\[
 2\varepsilon c_n
 =\frac{3n-4}{2}c_{n-1}
  -\sum_{j=1}^{n-1}c_jc_{n-j},
 \qquad n\ge1.
 \tag{8}
\]

$\varepsilon=1$ のとき
$n=1$ で $c_1=-1/4$ である．
式 (8) から帰納的に $c_n<0$（$n\ge1$）が従う．
実際 $n\ge2$ では第一項は負で，
和の各項は正なので右辺は負である．
さらに
\[
 |c_n|\ge\frac{3n-4}{4}|c_{n-1}|
 \qquad(n\ge2)
\]
を得る．
$\varepsilon=-1$ のときも
式 (8) から $c_1=-1/4$ となり，
帰納的に $(-1)^nc_n>0$（$n\ge1$）である．
この場合も第一項と和の負号が同じ符号を持つので
\[
 |c_n|\ge\frac{3n-4}{4}|c_{n-1}|
 \qquad(n\ge2)
\]
が成立する．
従ってどちらの場合も
$|c_n|$ は $n!$ と同程度以上の速さで増加し，
級数 $\sum c_nz^n$ の収束半径は零である．
しかし代数関数のPuiseux級数はある穿孔円板で
収束するはずである．これは矛盾する．
よって代数的なRiccati解は存在しない．
\end{proof}

\begin{theorem}[Airy方程式の非可積分性]
\label{thm:airy}
\[
 y''=xy,\qquad K=\mathbb C(x),\quad
 \partial=\frac d{dx}
\]
の基本解を含むLiouvillian拡大は存在しない．
従ってその微分Galois群の単位成分は非可解である．
\end{theorem}
\begin{proof}
もしLiouvillianな基本解系が存在すれば，
定理\ref{thm:riccati-criterion}から
Riccati方程式 $u'+u^2=x$ が
$\mathbb C(x)$ 上代数的な解を持つはずである．
これは補題\ref{lem:airy-riccati}に反する．
従ってLiouvillianな基本解系は存在しない．
Kolchinの定理からGalois群の単位成分も非可解である．
\end{proof}

\begin{remark}
代数方程式のGalois理論では根号可解性が有限群の可解性に
対応する．微分Galois理論では，積分・指数積分・代数拡大を
含むLiouvillian可積分性が，
線形代数群の\emph{単位成分}の可解性に対応する．
有限な非可解商群が存在しても単位成分は自明であり，
代数的拡大を許すLiouvillian可積分性を妨げないことに注意する．
\end{remark}

\section{到達課題の完全解答}
\noindent\textbf{一次方程式のLiouvillian解．}
\begin{proposition}
$a\in\mathbb C(x)$ のとき
$y'=a(x)y$ の非零解は，
微分体を必要に応じてLiouvillian拡大することで
\[
y=\exp\left(\int a(x)\,dx\right)
\]
と表せる．
\end{proposition}
\begin{proof}
$a$ の原始関数 $F$ を添加した微分体を考える．
次に $z'=F'z=az$ を満たす指数型の元 $z$ を添加する．
この二段階はそれぞれ原始関数の添加と
指数型の添加であり，Liouvillian拡大の操作である．
$z$ は与えられた方程式の非零解となる．
\end{proof}
\begin{proposition}
$y'=(1/x)y$（$x\ne0$）は有理解 $y=x$ を持つ．
\end{proposition}
\begin{proof}
$y=x$ と置けば $y'=1$，
$(1/x)y=1$ で一致する．
従ってこの方程式は基礎微分体 $\mathbb C(x)$
の内部にすでに非零解を持つ．
\end{proof}
\noindent\textbf{到達確認．}
$y'=3xy$ の解を求めよ．
\par\noindent\textbf{解答．}
$\int3x\,dx=3x^2/2$ だから
$y=C\exp(3x^2/2)$ である．
これは指数型の添加によるLiouvillian解である．
\end{document}